% 综述主文档模板
% 基于UAV_Delay_Survey_v3成功经验

\documentclass[12pt,a4paper]{ctexart}

% 基础包
\usepackage{geometry}
\usepackage{graphicx}
\usepackage{amsmath,amssymb,amsfonts}
\usepackage{booktabs}
\usepackage{multirow}
\usepackage{hyperref}
\usepackage{bookmark}
\usepackage{listings}
\usepackage{xcolor}
\usepackage{tikz}
\usetikzlibrary{shapes.geometric, arrows.meta, positioning, calc}

% 页面设置
\geometry{left=2.5cm,right=2.5cm,top=2.5cm,bottom=2.5cm}

% 标题设置
\title{\textbf{{TITLE}}}
\author{{AUTHOR}}
\date{{DATE}}

% TikZ样式定义
\tikzset{
    block/.style={rectangle, rounded corners, draw=black!70, fill=blue!10, 
                  minimum width=2.5cm, minimum height=1cm, text centered},
    arrow/.style={-{Stealth[length=2.5mm]}, thick, black!70},
    label/.style={font=\small\bfseries}
}

\begin{document}

\maketitle

% 摘要
\begin{abstract}
{ABSTRACT}

\textbf{关键词：}{KEYWORDS}
\end{abstract}

% 目录
\tableofcontents
\newpage

% 第1章 引言
\section{引言}
\subsection{研究背景与意义}
{CH1_1}

\subsection{问题定义与挑战}
{CH1_2}

\subsection{研究现状概述}
{CH1_3}

\subsection{本文组织结构}
{CH1_4}

% 第2章 基础理论
\section{基础理论}
\subsection{核心概念定义}
{CH2_1}

\subsection{数学建模}
{CH2_2}

\subsection{经典方法回顾}
{CH2_3}

% 第3章 方法A
\section{方法A}
\subsection{子类A1}
{CH3_1}

\subsection{子类A2}
{CH3_2}

\subsection{子类A3}
{CH3_3}

\subsection{方法对比与局限性}
{CH3_4}

% 第4章 方法B
\section{方法B}
\subsection{子类B1}
{CH4_1}

\subsection{子类B2}
{CH4_2}

\subsection{子类B3}
{CH4_3}

\subsection{方法对比与局限性}
{CH4_4}

% 第5章 方法C
\section{方法C}
\subsection{子类C1}
{CH5_1}

\subsection{子类C2}
{CH5_2}

\subsection{子类C3}
{CH5_3}

\subsection{方法对比与局限性}
{CH5_4}

% 第6章 实验与应用
\section{实验与应用}
\subsection{数据集与评测指标}
{CH6_1}

\subsection{性能对比分析}
{CH6_2}

\subsection{应用案例}
{CH6_3}

% 第7章 未来方向
\section{未来研究方向}
\subsection{技术挑战}
{CH7_1}

\subsection{研究趋势}
{CH7_2}

\subsection{发展展望}
{CH7_3}

% 参考文献
\begin{thebibliography}{99}
{BIBLIOGRAPHY}
\end{thebibliography}

\end{document}
